% Insert after the full finite coupling fibre theorem and reconstruction.
\subsection{A proper algebraic exceptional set for the full spatial projection}
\label{q444:generic-projection-section}

We now retain all sixteen spatial coordinates simultaneously. The new
statement is about the entire spatial projection, not only a selected
block class or one value of the coupling. The five original equations and
their prescribed rotating ansatz remain unchanged.

Use the exact fixed matrix $U$ of
\eqref{q444:complex-eight-coordinates}. For each $j=1,2$, write
$\widetilde Z_j=U^*Z_jU$ using the eight real coordinates
$(a_j,h_j,\operatorname{Re}v_{j1},\operatorname{Im}v_{j1},
\operatorname{Re}v_{j2},\operatorname{Im}v_{j2},
\operatorname{Re}z_j,\operatorname{Im}z_j)$, with $b_j=(h_j-a_j)/2$.
This is an invertible linear coordinate map on the original spatial
carrier. Denote the resulting sixteen-coordinate vector by $\xi$.
No coordinate is held at zero in the theorem below.

\begin{lemma}[Polynomial diagonal residuals with retained denominators]
\label{q444:generic-polynomial-diagonals}
Let $\eta=\langle J,\mathcal KJ\rangle$ be the original nonnegative
polynomial in $\xi$. On $\eta>0$ define
\[
 F_\xi(t)=\eta^{16}(U^*P_H(t)U)_{11},\qquad
 G_\xi(t)=\eta^{16}(U^*P_H(t)U)_{22}.
\]
Both extend as polynomials in $(\xi,t)$ with real algebraic
coefficients, and their degree in $t$ is at most twelve.
On the stated domain the inverse formulas are exactly
$(U^*P_HU)_{11}=F_\xi/\eta^{16}$ and
$(U^*P_HU)_{22}=G_\xi/\eta^{16}$.
\end{lemma}
\begin{proof}
All real coordinate components of $\mathcal K,w,\eta,\beta,\Pi\mathscr B$
are polynomials in $\xi$ with fixed real algebraic coefficients; the
projection denominator $\langle J,J\rangle=3/4$ is a fixed nonzero
constant. Choose once and for all a fixed basis of $W$ with real
algebraic coordinates in the displayed diagonal-frame carrier.
The entries of $\mathscr M=t\mathcal I_7+2\mathcal S_7$ have
denominator dividing $\eta$, and those of $f_7$ have the same property.
The determinant expansion shows that $\eta^7D_7$ is polynomial.
Each adjugate entry is a six-by-six minor, so $\eta^7V_7$ is
polynomial as well. The full formula
\eqref{q444:finite-coupling-N} therefore makes $\eta^8N_7$
polynomial. Complex conjugation here conjugates fixed coefficients and
the displayed complex coordinate combinations, while every component of
$\xi$ and $t$ is real. Thus $\eta^{14}Q_7$ is polynomial.
The full formula for $P_H$ now proves polynomial extension after the
displayed multiplication by $\eta^{16}$. No value at $\eta=0$
is treated as an inverse Gauss solution.

To prove the sharper degree bound, work through the exact conjugation
by $U^*$, so $H_d=\operatorname{diag}(0,r,-r)$ and
$J_d=iH_d$ are both diagonal. Write $N_7=V_7+s_7J$ as before.
Since $[J,H]=0$,
\[
 [N_7,[N_7,H]]=[V_7,[V_7,H]]+s_7[J,[V_7,H]].
\]
For every matrix $Y$, $([J_d,Y])_{jj}=0$ follows by subtracting
$(J_d)_{jj}Y_{jj}$ from itself. Consequently the second term contributes
zero to each diagonal entry after conjugation. The first term has
degree at most twelve since $\deg_t V_7\le6$. The other term in
$P_H$ is $Q_7$ times a matrix independent of $t$, and
$\deg_t Q_7\le12$. Multiplication by $\eta^{16}$ does not
change the degree in $t$. This proves every assertion.
\end{proof}

Define a fixed-size polynomial matrix $\mathscr S(\xi)$ by the
linear map
\begin{equation}
 \mathbb R[t]_{<12}\oplus\mathbb R[t]_{<12}
 \longrightarrow\mathbb R[t]_{<24},\qquad
 (A,B)\longmapsto AF_\xi+BG_\xi,
 \qquad \mathscr R(\xi)=\det\mathscr S(\xi).
 \label{q444:generic-sylvester-map}
\end{equation}
Use the increasing-power bases $1,t,\ldots$ and put the $A$ summand
before the $B$ summand. This specifies the determinant sign and its
twenty-four columns, including when either polynomial loses degree.
This is the classical Sylvester resultant map; see the Stacks Project,
\href{https://stacks.math.columbia.edu/tag/00UA}{Tag 00UA}.
The argument below proves the properties needed here directly, without
replacing either polynomial by a monic one.

\begin{theorem}[Generic exclusion in all sixteen spatial coordinates]
\label{q444:generic-projection-obstruction}
The explicitly defined polynomial $\mathscr R$ is not identically
zero. The spatial projection of every full solution of
\eqref{q444:cx:eq:phi}--\eqref{q444:cx:eq:sp2} is contained in
\[
 \{\xi\in\mathbb R^{16}:\mathscr R(\xi)=0\}.
\]
This exceptional set is closed, has empty interior, and has
sixteen-dimensional Lebesgue measure zero. Thus the complement is an
open dense set of unrestricted spatial pairs, each admitting no positive
coupling and no frequency sign solving the five retained equations.
This theorem does not decide whether any point of the exceptional set
admits a solution.
\end{theorem}
\begin{proof}
Any full solution has $[Z_1,Z_2]\ne0$, hence $\eta>0$ by
Lemma~\ref{q444:finite-coupling-schur}. It also has $Q_7(t)>0$ at
the original $t=g^2$, and the exact residual identity
\eqref{q444:finite-coupling-exact-residual-map} makes the entire
$P_H(t)$ zero. Therefore $F_\xi(t)=G_\xi(t)=0$. Evaluation at
that real $t$ is a nonzero linear functional on
$\mathbb R[t]_{<24}$ (it sends $1$ to $1$), and it annihilates
the image of \eqref{q444:generic-sylvester-map}. That image cannot
be all of the codomain. The determinant is zero. This proof includes
every degree-drop case and does not use a converse resultant criterion
outside the original denominator domain.

Here is an exact witness that the determinant polynomial is nonzero.
Write
\[
 M(a,h,v_1,v_2,z)=
 \begin{pmatrix}
 ia&-\bar v_1&-\bar v_2\\
 v_1&i(h-a)/2&z\\
 v_2&-\bar z&-i(a+h)/2
 \end{pmatrix}
\]
and set, in the displayed diagonal coordinates,
\[
 \widetilde Z_1=M(1,2,1+i,-1+2i,2-i),\qquad
 \widetilde Z_2=M(-2,1,2+3i,1-i,-1+2i).
\]
The original pair is $Z_j=U\widetilde Z_jU^*$, not a replacement
of the prescribed $H,q,F$. Every one of its sixteen displayed real
spatial coordinates is nonzero. Substitution in the original definitions
of $D_7,V_7,N_7,Q_7,P_H$ gives two diagonal polynomials of degree
twelve in $\mathbb Q(\sqrt6)[t]$. Their coprimality can be verified
by the following small certificate, with all coefficient lists in
descending powers. Reduce their coefficients by the explicit map
$\sqrt6\mapsto39$ to $\mathbb F_{101}$, where $39^2=6$.
Each rational coefficient denominator is nonzero modulo $101$;
the exact coefficient file and replay test check them individually.
The two reduced polynomials $f,g$ and the Bezout coefficients $a,b$ are
\begin{align*}
 f &: (39,85,79,93,67,20,38,98,26,56,8,25,28),\\
 g &: (1,87,96,7,76,88,75,93,53,73,64,41,16),\\
 a &: (4,60,31,59,74,31,34,18,59,97,95,96),\\
 b &: (46,85,47,5,14,36,38,21,35,48,34,53).
\end{align*}
Multiplication gives $af+bg=1$ in $\mathbb F_{101}[t]$.
Both $f$ and $g$ have degree exactly twelve. If a pair
$(a_0,b_0)$ of degrees less than twelve were in the kernel of their
Sylvester map, coprimality would imply that $f$ divides $b_0$ and
$g$ divides $a_0$. The degree bounds force both to be zero. Thus that
finite-field map, between spaces of equal dimension twenty-four, has
nonzero determinant. Reduction commutes with the determinant polynomial.
It follows that the Sylvester determinant of the two original diagonal
polynomials is nonzero in $\mathbb Q(\sqrt6)$ as well. This uses
the localized coefficient ring with all its denominators inverted,
not a homomorphism from the whole field $\mathbb Q$ to a finite field.

For explicit verification at the witness, multiplication gives
\[
 [\widetilde Z_1,\widetilde Z_2]_{21}=9-\tfrac72 i\ne0,\qquad
 \eta=\tfrac34(7+15)+3(5+5)=\tfrac{93}2>0.
\]
Multiplying both diagonal polynomials by the same nonzero scalar
$\eta^{16}$ multiplies each of the twenty-four columns by
$\eta^{16}$, and hence multiplies their determinant by
$\eta^{384}\ne0$. Consequently $\mathscr R$ is nonzero at
this exact witness and cannot be the zero polynomial. The finite-field
calculation is solely a certificate of this real algebraic statement;
it does not change the physical carrier or any coupling.

As a separate certificate, write $j_0=\langle J,J\rangle=3/4$.
In $W\oplus\mathbb RJ$, without changing the length of $J$, the
full Gauss endomorphism has block matrix
\[
 L=\begin{pmatrix}
 t\mathcal I_7+2\Pi\mathcal K|_W&2w\\
 2\langle w,\cdot\rangle/j_0&2\eta/j_0
 \end{pmatrix}.
\]
Subtracting $j_0w/\eta$ times its last row from the upper block rows
preserves the determinant and makes the upper-right block zero; its
upper-left block becomes $\mathscr M(t)$. Thus
$\det L=(2\eta/j_0)D_7$ exactly.
The independent replay reconstructs both original diagonal polynomials
from that matrix modulo $101$. It obtains Sylvester determinant $40$;
after the retained $\eta^{16}$ multiplication it obtains
$40\eta^{384}=60\ne0$ in $\mathbb F_{101}$. These residues refer
to the exact basis order in \eqref{q444:generic-sylvester-map}.

For completeness, a nonzero real polynomial on $\mathbb R^n$ has a
measure-zero zero set by induction on $n$. For $n=1$ it has finitely
many roots. For $n>1$, write it as a polynomial in the last variable
whose coefficients are polynomials in the first $n-1$ variables.
Choose one nonzero coefficient. The set where all coefficients vanish
is contained in the measure-zero zero set of that coefficient, by
induction. At every other base point the last-variable polynomial has
only finitely many roots. Applying Fubini on each bounded box shows
that the zero set in that box has measure zero; countably many boxes
cover $\mathbb R^n$. Its complement is open because the polynomial
is continuous. Every nonempty open ball has positive measure, so no
ball is contained in the zero set, and the complement is dense.
This proves the assertions at $n=16$.
\end{proof}

The exact large-coefficient identity over $\mathbb Q(\sqrt6)$ and
the small finite-field identity are retained separately in
\path{PAIR_BEZOUT_CERTIFICATE.json} and
\path{PAIR_MODULAR_CERTIFICATE.json}. The replay
\path{check_pair_bezout.py} constructs the original full pair and every
Gauss and scalar numerator, then checks the identities; the integer-only
certificate replay checks the displayed finite-field products separately.
Neither a finite failed search nor two failed test couplings proves this
theorem: the proof uses an explicit nonzero polynomial vanishing on
every possible spatial projection, with its nonzeroness certified at
the literal witness.
